% =====================================================================
%  1.0 EXECUTIVE SUMMARY  (5 pts)   P1.2, 31 Oct 2026
%  Rubric (5.1.4):  [x] Objective statement
%                   [x] Planned approach to achieve all objectives
%                   [x] Main points from every subsequent section (budget and funding included, P1.2)
% =====================================================================
\section{Executive Summary}
% P1.2: the side panel became a full-width readout strip so page 1 no longer grows with the panel
% (and stays fixed when the team's drawings arrive; the GA plate in 3.4 carries them).
\newcommand\gl[3]{\parbox[t]{0.163\linewidth}{\raggedright
  {\dbflabelfont\fontsize{7}{8.4}\selectfont\color{dbfFog}\MakeUppercase{#1}}\par\vspace{2pt}%
  {\dbftablefont\bfseries\fontsize{12.5}{14.5}\selectfont\color{white}#2}\hspace{0.2em}{\sffamily\fontsize{7.6}{9}\selectfont\color{dbfFog}#3}}}
\textbf{Objective.} Machine Masters will design, build and fly a hand-flown electric aircraft that delivers the sensor in its container (Mission~2), tows it through its light sequence and recovers it (Mission~3), and protects it in the Ground Mission drop, to post the highest total score at Tucson in April 2027.

\textbf{Approach.} The rules became 21 subsystem requirements. A mission model calibrated on the 2022--26 fly-offs flew 162 candidate aircraft through every mission, failed attempts included, against simulated rival fields; completion pays 4.5 of 7.5 points, so reliability carries \tCompletion\,\% of the trade weight. Sensor and container are built for 3--8~kg, with the maximum declared at inspection ($+$\tDeclGain{} points over fixing it at design).

\textbf{Design.} A \nSpan~m, \nWingArea~m² high wing with Fowler flaps ($C_{L,\max}$ \nCLmaxLand), conventional tail, fixed tricycle gear and one 8S tractor stalls at \nStallMtwo~m/s at its \nGrossMtwo~kg Mission~2 mass. It waives Mission~1, carries one \nMtwoPayload~kg sensor container, tows \nMthreeMass~kg and drops the declared sensor from \nDrop~m: \sTotal{} of 7.5 with every mission completed and the field's highest mission total in \sPtopPct\,\% of simulated contests, failed attempts counted (first iteration \sPtopCtlPct\,\%); a slotted 1.2~m² wing is the fallback.

\textbf{Execution.} Six sub-teams build two airframes of laser-cut ply and balsa on carbon tubes, retire the top risks in 14 bench and flight tests by 15~January, fly the prototype at least 20 times and freeze the design on 19~February; the \$\bTotalK{}k budget, travel included, is fully funded by the university.

\vspace{6pt}\noindent\begin{dbfinkpanel}[Design at a glance / Variant A]
\gl{Wingspan}{\nSpan}{m}\hfill\gl{Wing area}{\nWingArea}{m²}\hfill\gl{M2 gross mass}{\nGrossMtwo}{kg}\hfill\gl{Stall, M2}{\nStallMtwo}{m/s}\hfill\gl{Score, if done}{\sTotal}{/ 7.5}\hfill\gl{P(top of field)}{\sPtop}{}
\end{dbfinkpanel}
